% !PRINTED 474
\weberpage{474}
%<<CONT
\gesperrt{Kronecker} recognised\footnote{\gesperrt{Kronecker}, \glqq Ueber Gleichungen 7\textsuperscript{ten} Grades\grqq{}. Monatsberichte der Berliner Akademie 1858.}, and its closer investigation is of great importance for the general theory of the equation of degree 7.

We can obtain this permutation group by seeking out the permutations which are produced by the substitutions of the group $G_{168}$ among the quantities $z_0, z_1, z_2, z_3, z_4, z_5, z_6$ (\S.\,125). Here it is then to be observed, with regard to composition, that if the two linear substitutions $\xi_1, \xi_2$ effect the permutations $\pi_1, \pi_2$, the composed permutation $\pi_1 \pi_2$ is produced by $\xi_2 \xi_1$. For according to the definition one obtains $\xi_2 \xi_1 (x)$ by applying to the variables $(x)$ first the substitution $\xi_1$ and to the result the substitution $\xi_2$. Likewise $\pi_1 \pi_2$ denotes the permutation which results when one applies to the seven figures first $\pi_1$ and then $\pi_2$. We thus obtain, corresponding to the group $G_{168}$, a permutation group of degree 168, which we denote by $P_{168}$, and these two groups are isomorphic.

Since we have recognised $\tau$ and $\omega$ as generating elements of the group $G_{168}$ (\S.\,72, I.), it suffices to determine the permutations corresponding to these two substitutions in order to derive from them the whole group $P_{168}$.

Now from the substitution $\tau$ there arises the cyclic permutation of the seven figures

\[
(\tau) \qquad (0, 1, 2, 3, 4, 5, 6)
\]

and the influence of $\omega$ follows from the remark that $z_0$ remains unaltered by the substitutions of the octahedral group $\chi^\lambda \omega^u \Theta^v$. Hence, to obtain the alteration of $z_r$ by $\omega$, we have only to ascertain the influence of $\omega \tau^r$ on $z_0$. But this influence follows immediately from the formulae \S.\,72, (15), according to which, e.g., $\omega \tau = \tau^2 \chi^2 \omega \Theta^2$, so that $z_1$ passes into $z_2$, and so on. Accordingly, to the substitution $\omega$ corresponds the permutation

\[
(\omega) \qquad \begin{pmatrix} 0, & 1, & 2, & 3, & 4, & 5, & 6 \\ 0, & 2, & 1, & 5, & 4, & 3, & 6 \end{pmatrix} = (1, 2)\,(3, 5).
\]

One can represent the group $P_{168}$ by the congruence group of ter-
%CONT>>
